\documentclass[12pt,a4paper]{article}
\usepackage[utf8]{inputenc}
\usepackage[T1]{fontenc}
\usepackage{amsmath,amssymb}
\usepackage{graphicx}
\usepackage{listings}
\usepackage{xcolor}
\usepackage{colortbl}
\usepackage{hyperref}
\usepackage{geometry}
\usepackage{tikz}
\usepackage[most]{tcolorbox}
\usepackage{needspace}
\usepackage{mdframed}
\usetikzlibrary{trees,positioning,arrows,shapes}
\geometry{margin=2.5cm}

% Définition des couleurs personnalisées
\definecolor{titrebleu}{RGB}{0,102,204}
\definecolor{defvert}{RGB}{0,153,76}
\definecolor{exempleorange}{RGB}{255,102,0}
\definecolor{algoviolet}{RGB}{153,0,153}
\definecolor{importantrouge}{RGB}{204,0,0}
\definecolor{fondgris}{RGB}{245,245,245}
\definecolor{fondjaune}{RGB}{255,250,205}
\definecolor{fondvert}{RGB}{230,255,230}
\definecolor{fondbleu}{RGB}{230,240,255}
\definecolor{fondrose}{RGB}{255,230,240}

% Configuration des listings
\lstset{
    basicstyle=\ttfamily\small,
    breaklines=true,
    frame=single,
    numbers=left,
    numberstyle=\tiny\color{gray},
    showstringspaces=false,
    commentstyle=\color{green!60!black},
    keywordstyle=\color{blue!80!black}\bfseries,
    stringstyle=\color{red!70!black},
    backgroundcolor=\color{fondgris},
    rulecolor=\color{titrebleu}
}

% Boîtes colorées - AVEC BREAKABLE
\newtcolorbox{definition}{
    colback=fondvert,
    colframe=defvert,
    fonttitle=\bfseries,
    title=Définition,
    arc=3mm,
    breakable
}

\newtcolorbox{exemple}{
    colback=fondjaune,
    colframe=exempleorange,
    fonttitle=\bfseries,
    title=Exemple,
    arc=3mm,
    breakable
}

\newtcolorbox{important}{
    colback=white,
    colframe=importantrouge,
    fonttitle=\bfseries,
    title=Important,
    arc=3mm,
    breakable
}

\newtcolorbox{algorithme}{
    colback=fondbleu,
    colframe=algoviolet,
    fonttitle=\bfseries,
    title=Algorithme,
    arc=3mm,
    breakable
}

\newtcolorbox{notecours}{
    colback=fondrose,
    colframe=importantrouge,
    fonttitle=\bfseries,
    title=Note de Cours,
    arc=3mm,
    breakable
}

% Style des sections
\usepackage{titlesec}
\titleformat{\section}
{\color{titrebleu}\normalfont\Large\bfseries}
{\color{titrebleu}\thesection}{1em}{}

\titleformat{\subsection}
{\color{defvert}\normalfont\large\bfseries}
{\color{defvert}\thesubsection}{1em}{}

\title{\textbf{\Huge\color{titrebleu}Types Abstraits de Données}\\[0.5em]
\large\color{defvert}Structures de Données et Implémentations}
\author{\large Cours détaillé avec exemples et notes complémentaires}
\date{\large Version 2025 - Édition Complète}

\begin{document}

\maketitle
\tableofcontents
\newpage

\section{Introduction}

Ce document présente les structures de données fondamentales et les types abstraits de données (TAD).

\subsection{Distinction Structure vs Type Abstrait}

\begin{definition}
\textbf{Structure de données :} Implémentation concrète en mémoire (tableaux, listes chaînées, arbres).

\textbf{Type Abstrait de Données (TAD) :} Spécification abstraite des opérations, indépendante de l'implémentation.
\end{definition}

\begin{notecours}
\textbf{Interface ici, c'est la liste de fonction valides.}

Un TAD définit clairement quelles opérations sont possibles sans se préoccuper de comment elles sont implémentées en interne.
\end{notecours}

\begin{exemple}
\begin{itemize}
    \item \textbf{TAD :} \textit{Pile} (opérations : empiler, dépiler)
    \item \textbf{Structures possibles :} tableau, liste chaînée
\end{itemize}
\end{exemple}

\newpage
\section{Structures de Données Fondamentales}

\subsection{La Pile (Stack)}

\subsubsection{Définition}

\begin{definition}
Une \textbf{pile} est une structure de données \textcolor{importantrouge}{\textbf{LIFO}} (Last In, First Out) : le dernier élément inséré sera le premier à être traité.

\textbf{Principe :} Comme une pile d'assiettes - on ajoute et retire toujours par le dessus.
\end{definition}

\begin{notecours}
\textbf{Points importants sur la pile :}
\begin{itemize}
    \item En \textbf{violet} = \textbf{état temporaire actuel} pour l'exemple
    \item \textbf{taille\_pile = 5} dans l'exemple illustré
    \item \textbf{L'index commence toujours par 1} dans ce cours (il n'y a pas d'index 0!)
\end{itemize}

\textbf{Notes manuscrites :}
\begin{itemize}
    \item \textbf{Empiler c'est un procédure}
    \item \textbf{Dépiler c'est une fonction} qui change l'état du mémoire et retourne la valeur à la tête du pile
\end{itemize}
\end{notecours}

\subsubsection{Opérations}

\begin{important}
\begin{itemize}
    \item \textbf{empiler(x)} ou \textit{push(x)} : ajoute un élément au sommet
    \item \textbf{dépiler()} ou \textit{pop()} : retire et retourne l'élément du sommet
    \item \textbf{pileVide?()} ou \textit{isEmpty()} : teste si la pile est vide
    \item \textbf{sommet()} ou \textit{peek()} : consulte le sommet sans retirer
\end{itemize}
\end{important}

\subsubsection{Implémentation par tableau}

\begin{lstlisting}[language=Pascal]
Type Pile = Record
    tab : Array[1..MAX] of Integer;
    taille : Integer;      // nombre d'elements
    tailleMax : Integer;   // capacite maximale
End;
\end{lstlisting}

\begin{exemple}
\textbf{Représentation graphique :}
\begin{verbatim}
Indices:   1   2   3   4   5   6   7   8   9  10
        +---+---+---+---+---+---+---+---+---+---+
tab[]   | 1 | 4 |12 | 8 | 9 |14 |20 | 5 | 6 | 2 |
        +---+---+---+---+---+---+---+---+---+---+
                                            ^
                                        taille=10
\end{verbatim}
\end{exemple}

\subsubsection{Algorithmes}

\begin{algorithme}
\textbf{Empiler :}
\begin{verbatim}
Procédure empiler(p : Pile, x : Integer)
Début
    Si p.taille >= p.tailleMax Alors
        Erreur("Pile pleine")
    Sinon
        p.taille <- p.taille + 1
        p.tab[p.taille] <- x
    FinSi
Fin
\end{verbatim}
\end{algorithme}

\begin{algorithme}
\textbf{Dépiler :}
\begin{verbatim}
Fonction dépiler(p : Pile) : Integer
Début
    Si p.taille = 0 Alors
        Erreur("Pile vide")
    Sinon
        x <- p.tab[p.taille]
        p.taille <- p.taille - 1
        Retourner x
    FinSi
Fin
\end{verbatim}
\end{algorithme}

\subsubsection{Exemple d'utilisation}

\begin{exemple}
\begin{lstlisting}[language=Python]
# Creation d'une pile
pile = []

# Empilement
pile.append(10)  # [10]
pile.append(20)  # [10, 20]
pile.append(30)  # [10, 20, 30]

# Depilement
x = pile.pop()   # x = 30, pile = [10, 20]
y = pile.pop()   # y = 20, pile = [10]

# Application : verification de parentheses
def parentheses_equilibrees(expression):
    pile = []
    for char in expression:
        if char == '(':
            pile.append(char)
        elif char == ')':
            if not pile:
                return False
            pile.pop()
    return len(pile) == 0

# Test
print(parentheses_equilibrees("((a+b)*(c-d))"))  # True
print(parentheses_equilibrees("((a+b)"))          # False
\end{lstlisting}
\end{exemple}

\subsubsection{Complexités}

\begin{center}
\colorbox{fondbleu}{
\begin{tabular}{|l|c|}
\hline
\rowcolor{titrebleu!30}
\textbf{Opération} & \textbf{Complexité} \\
\hline
empiler & \textcolor{defvert}{\textbf{O(1)}} \\
dépiler & \textcolor{defvert}{\textbf{O(1)}} \\
pileVide? & \textcolor{defvert}{\textbf{O(1)}} \\
sommet & \textcolor{defvert}{\textbf{O(1)}} \\
\hline
\end{tabular}}
\end{center}

\newpage
\subsection{La File (Queue)}

\subsubsection{Définition}

\begin{definition}
Une \textbf{file} est une structure de données \textcolor{importantrouge}{\textbf{FIFO}} (First In, First Out) : le premier élément inséré sera le premier à être traité.

\textbf{Principe :} Comme une file d'attente - on entre par un bout, on sort par l'autre.
\end{definition}

\subsubsection{Opérations}

\begin{important}
\begin{itemize}
    \item \textbf{enfiler(x)} ou \textit{enqueue(x)} : ajoute un élément en fin de file
    \item \textbf{défiler()} ou \textit{dequeue()} : retire et retourne l'élément en début
    \item \textbf{fileVide?()} ou \textit{isEmpty()} : teste si la file est vide
\end{itemize}
\end{important}

\subsubsection{Implémentation par tableau circulaire}

\begin{lstlisting}[language=Pascal]
Type File = Record
    tab : Array[0..MAX-1] of Integer;
    sortie : Integer;   // indice de sortie
    entree : Integer;   // indice d'entree
    taille : Integer;   // nombre d'elements
End;
\end{lstlisting}

\begin{exemple}
\textbf{Représentation graphique :}
\begin{verbatim}
        sortie              entree
          |                   |
          v                   v
    +---+---+---+---+---+---+---+---+---+---+
    | 8 | 9 |14 |20 | 5 | 6 | 2 | 3 |   |   |
    +---+---+---+---+---+---+---+---+---+---+
      0   1   2   3   4   5   6   7   8   9
      
    taille = 8
\end{verbatim}
\end{exemple}

\subsubsection{Algorithmes}

\begin{algorithme}
\textbf{Enfiler :}
\begin{verbatim}
Procédure enfiler(f : File, x : Integer)
Début
    Si f.taille = MAX Alors
        Erreur("File pleine")
    Sinon
        f.tab[f.entree] <- x
        f.entree <- (f.entree + 1) MOD MAX
        f.taille <- f.taille + 1
    FinSi
Fin
\end{verbatim}
\end{algorithme}

\begin{algorithme}
\textbf{Défiler :}
\begin{verbatim}
Fonction défiler(f : File) : Integer
Début
    Si f.taille = 0 Alors
        Erreur("File vide")
    Sinon
        x <- f.tab[f.sortie]
        f.sortie <- (f.sortie + 1) MOD MAX
        f.taille <- f.taille - 1
        Retourner x
    FinSi
Fin
\end{verbatim}
\end{algorithme}

\subsubsection{Implémentation par double pile}

\begin{important}
Une file peut aussi être implémentée avec deux piles :
\begin{itemize}
    \item \textbf{pile\_entree} : pour les enfilements
    \item \textbf{pile\_sortie} : pour les défilements
\end{itemize}
\end{important}

\begin{notecours}
\textbf{Observations importantes (notes manuscrites) :}
\begin{itemize}
    \item \textbf{EN BLEU + BLANCHE = NON UTILISÉ} (zones vides)
    \item \textbf{VIOLET = UTILISÉ} (zones contenant des données)
    \item L'implémentation abstraite utilise une \textbf{liste chaînée} $\rightarrow$ code facile
    \item Le code à faire en copiant sur la pile
    \item Pour file vide $\rightarrow$ p1 vide ET p2 vide
    \item \textbf{Pour file vide --> p1 vide p2 vide}
\end{itemize}
\end{notecours}

\begin{exemple}
\begin{lstlisting}[language=Python]
class FileDoublePile:
    def __init__(self):
        self.pile_entree = []
        self.pile_sortie = []
    
    def enfiler(self, x):
        self.pile_entree.append(x)
    
    def defiler(self):
        # Si pile_sortie vide, transferer pile_entree
        if not self.pile_sortie:
            while self.pile_entree:
                self.pile_sortie.append(
                    self.pile_entree.pop()
                )
        
        if not self.pile_sortie:
            raise Exception("File vide")
        
        return self.pile_sortie.pop()

# Exemple d'utilisation
f = FileDoublePile()
f.enfiler(1)
f.enfiler(2)
f.enfiler(3)
print(f.defiler())  # 1
print(f.defiler())  # 2
\end{lstlisting}
\end{exemple}

\newpage
\subsection{Arbre Binaire de Recherche (ABR)}

\subsubsection{Définition}

\begin{definition}
Un \textbf{arbre binaire de recherche} est un arbre binaire tel que pour tout nœud :
\begin{itemize}
    \item Les valeurs du sous-arbre \textcolor{defvert}{\textbf{gauche}} sont $\leq$ à la valeur du nœud
    \item Les valeurs du sous-arbre \textcolor{importantrouge}{\textbf{droit}} sont $>$ à la valeur du nœud
\end{itemize}
\end{definition}

\begin{notecours}
\textbf{Point crucial (notes manuscrites) :}

\textbf{Arbres binaires sont pas équitable} en termes de capacité (log(n) et n).

La complexité est en O($\log n$) pour un arbre équilibré, mais peut dégénérer en O(n) si l'arbre est déséquilibré (forme de liste).
\end{notecours}

\subsubsection{Structure}

\begin{lstlisting}[language=Pascal]
Type ABR = ^Noeud;
     Noeud = Record
         cle : Integer;
         sAG : ABR;      // sous-arbre gauche
         sAD : ABR;      // sous-arbre droit
         pere : ABR;     // pointeur vers le pere
     End;
\end{lstlisting}

\begin{exemple}
\textbf{Exemple d'ABR :}
\begin{verbatim}
           50
          /  \
        30    70
       /  \   / \
      20  40 60  80
     /              \
    10              90
\end{verbatim}
\end{exemple}

\subsubsection{Opérations principales}

\begin{algorithme}
\textbf{Recherche :}
\begin{verbatim}
Fonction rechercher(abr : ABR, x : Integer) : ABR
Début
    Si abr = NIL OU abr.cle = x Alors
        Retourner abr
    FinSi
    Si x < abr.cle Alors
        Retourner rechercher(abr.sAG, x)
    Sinon
        Retourner rechercher(abr.sAD, x)
    FinSi
Fin
\end{verbatim}
\end{algorithme}

\begin{exemple}
\textbf{Recherche de 40 :}
\begin{verbatim}
50 (40 < 50, aller à gauche)
  30 (40 > 30, aller à droite)
       40 (trouvé!)
\end{verbatim}
\end{exemple}

\begin{algorithme}
\textbf{Insertion :}
\begin{verbatim}
Procédure inserer(abr : ABR, x : Integer)
Début
    Si abr = NIL Alors
        abr <- NouveauNoeud(x)
    Sinon Si x < abr.cle Alors
        inserer(abr.sAG, x)
    Sinon
        inserer(abr.sAD, x)
    FinSi
Fin
\end{verbatim}
\end{algorithme}

\begin{exemple}
\textbf{Insertion de 65 :}
\begin{verbatim}
50 -> 70 (65 < 70, gauche)
  -> 60 (65 > 60, droite)
    -> insérer 65
\end{verbatim}
\end{exemple}

\subsubsection{Complexités}

\begin{center}
\colorbox{fondbleu}{
\begin{tabular}{|l|c|c|}
\hline
\rowcolor{titrebleu!30}
\textbf{Opération} & \textbf{Moyenne} & \textbf{Pire cas} \\
\hline
Recherche & \textcolor{defvert}{O($\log n$)} & \textcolor{importantrouge}{O(n)} \\
Insertion & \textcolor{defvert}{O($\log n$)} & \textcolor{importantrouge}{O(n)} \\
Suppression & \textcolor{defvert}{O($\log n$)} & \textcolor{importantrouge}{O(n)} \\
Minimum/Maximum & \textcolor{defvert}{O($\log n$)} & \textcolor{importantrouge}{O(n)} \\
\hline
\end{tabular}}
\end{center}

\newpage
\subsection{Tas (Heap)}

\subsubsection{Définition}

\begin{definition}
Un \textbf{tas} est un arbre binaire presque complet tel que pour tout nœud (sauf la racine), la valeur du père est supérieure ou égale à celle de ses fils.

\textbf{Tas max :} père $\geq$ fils (racine = maximum)

\textbf{Tas min :} père $\leq$ fils (racine = minimum)
\end{definition}

\begin{notecours}
\textbf{Propriété d'équilibre :}

Un tas est toujours \textbf{équilibré} car c'est un arbre binaire presque complet.

\textbf{Caractéristiques importantes (notes manuscrites) :}
\begin{enumerate}
    \item Un tas est une \textbf{interface avec 3 fonctions} principales : insertion, extraction du maximum, et consultation du maximum
    \item C'est une \textbf{structure presque pleine} : tout la structure sauf éventuellement le dernier niveau est complètement rempli de gauche à droite
    \item \textbf{C'est la structure de tas} elle-même (notion fondamentale à maîtriser)
    \item \textbf{Un tas est une interface avec 3 fonction}
    \item \textbf{Structure presque pleine tout la structure sauf et rempli de gauche a droit}
    \item \textbf{C'est la structutr de tas}
\end{enumerate}
\end{notecours}

\begin{exemple}
\textbf{Exemple de tas max :}
\begin{verbatim}
           90
          /  \
        80    70
       /  \   / \
      50  60 40  30
     / \
    10 20
\end{verbatim}
\end{exemple}

\subsubsection{Implémentation par tableau}

\begin{important}
Le tas est stocké dans un tableau où :
\begin{itemize}
    \item L'élément à l'indice i a :
    \item Père à l'indice $\lfloor i/2 \rfloor$
    \item Fils gauche à l'indice $2i$
    \item Fils droit à l'indice $2i + 1$
\end{itemize}
\end{important}

\begin{lstlisting}[language=Pascal]
Type Tas = Record
    tab : Array[1..MAX] of Integer;
    taille : Integer;
End;
\end{lstlisting}

\begin{exemple}
\textbf{Représentation du tas précédent :}
\begin{verbatim}
Indice:   1   2   3   4   5   6   7   8   9
       +---+---+---+---+---+---+---+---+---+
tab[]  |90 |80 |70 |50 |60 |40 |30 |10 |20 |
       +---+---+---+---+---+---+---+---+---+
\end{verbatim}
\end{exemple}

\begin{notecours}
\textbf{Opération d'insertion dans un tas (notes manuscrites) :}
\begin{itemize}
    \item Ici il faut ajouter le prochain élément en \textbf{fils gauche} du nœud num5
    \item Puis \textbf{monter si possible} (comparer avec le parent et échanger si nécessaire)
    \item Pour insérer dans un tas, on a une \textbf{complexité de hauteur} donc O($\log n$)
    \item \textbf{Ici il faur ajouter le prochain en fils gauche du noeud num5}
    \item \textbf{Puis monter si possible}
    \item \textbf{Pour insere tas on une complexite de hauteur donc log de n}
\end{itemize}
\end{notecours}

\subsubsection{Complexités}

\begin{center}
\colorbox{fondvert!20}{
\begin{tabular}{|l|c|}
\hline
\rowcolor{defvert!30}
\textbf{Opération} & \textbf{Complexité} \\
\hline
Insertion & \textcolor{algoviolet}{O($\log n$)} \\
Extraction max & \textcolor{algoviolet}{O($\log n$)} \\
Recherche max & \textcolor{defvert}{O(1)} \\
Construction & \textcolor{defvert}{O(n)} \\
\hline
\end{tabular}}
\end{center}

\begin{notecours}
\textbf{Notes manuscrites sur le tas :}
\begin{itemize}
    \item \textbf{On ajoute en bas a gauche puis on monte}
    \item \textbf{O(1) on prend le valeur au sommet}
    \item \textbf{Log(n) equilibre}
\end{itemize}
\end{notecours}

\newpage
\section{Types Abstraits de Données (TAD)}

\subsection{Définition générale}

\begin{definition}
Un \textbf{Type Abstrait de Données} est composé de :
\begin{enumerate}
    \item Un ensemble d'objets similaires
    \item Un ensemble d'opérations sur ces objets
\end{enumerate}

\textbf{Principe :} Séparation entre \textit{spécification} (quoi faire) et \textit{implémentation} (comment faire).
\end{definition}

\subsection{TAD Dictionnaire}

\begin{definition}
Un \textbf{dictionnaire} est un ensemble fini d'éléments issus d'un ensemble discret, muni d'une relation d'ordre.
\end{definition}

\begin{notecours}
\textbf{Notes importantes sur le dictionnaire (notes manuscrites) :}
\begin{itemize}
    \item \textbf{On a une dicotomie pour le log(n)} : à chaque itération, on diminue la taille par 2
    \item \textbf{Pour supprimer une valeur}, on prend le dernier élément et on le met à sa place
    \item Dans une \textbf{arbre de recherche}, pour supprimer : si c'est une feuille, c'est direct ; si c'est un nœud avec enfants, il faut le remplacer
    \item L'\textbf{équilibre} est de l'ordre de $\log n$ et c'est équilibré
    \item Pour le O(n) en tas, c'est parce qu'il y a des contraintes pour guider (gauche petit, droit grand), pas d'ordre strict
\end{itemize}
\end{notecours}

\subsubsection{Implémentations possibles}

\begin{center}
\colorbox{fondjaune}{
\begin{tabular}{|l|c|c|c|}
\hline
\rowcolor{exempleorange!30}
\textbf{Structure} & \textbf{Recherche} & \textbf{Insertion} & \textbf{Suppression} \\
\hline
Tableau non ordonné & O(n) & \textcolor{defvert}{O(1)} & \textcolor{defvert}{O(1)} \\
Liste non ordonnée & O(n) & \textcolor{defvert}{O(1)} & \textcolor{defvert}{O(1)} \\
Tableau ordonné & \textcolor{defvert}{O($\log n$)} & O(n) & O(n) \\
ABR & \textcolor{defvert}{O(h)} & \textcolor{defvert}{O(h)} & \textcolor{defvert}{O(h)} \\
\hline
\end{tabular}}
\end{center}

\subsection{TAD Ensemble Dynamique}

\begin{definition}
Un \textbf{ensemble dynamique} peut être modifié : insertion et suppression d'éléments.
\end{definition}

\begin{notecours}
\textbf{Points cruciaux (notes manuscrites) :}
\begin{itemize}
    \item Il faut savoir la \textbf{taille pour insérer en O(1)}
    \item Il s'intéresse \textbf{pas au tri} mais il faut connaître et le comprendre
\end{itemize}
\end{notecours}

\begin{important}
\textbf{Interface :} insertion, suppression, recherche d'un élément, de l'élément minimum, du maximum, du prédécesseur, du successeur.
\end{important}

\subsection{TAD File de Priorité}

\begin{definition}
Une \textbf{file de priorité} permet de gérer une liste d'éléments munis d'une clé, suivant un ordre de priorité défini sur les clés.

\textbf{Exemple d'application :} Gestion des processus dans un système d'exploitation.
\end{definition}

\begin{notecours}
\textbf{Notes manuscrites :}
\begin{itemize}
    \item \textbf{Idee que ca fonctinne comme un urgence avec gravite du cas}
    \item \textbf{File pile file de priorites sont des structures abstraites}
\end{itemize}
\end{notecours}

\subsubsection{Comparaison des implémentations}

\begin{center}
\colorbox{fondbleu}{
\begin{tabular}{|l|c|c|c|}
\hline
\rowcolor{titrebleu!30}
\textbf{Structure} & \textbf{Insérer} & \textbf{Maximum} & \textbf{ExtraireMax} \\
\hline
Tableau non ordonné & \textcolor{defvert}{O(1)} & O(n) & O(n) \\
Liste non ordonnée & \textcolor{defvert}{O(1)} & O(n) & O(n) \\
Tableau ordonné & O(n) & \textcolor{defvert}{O(1)} & O(n) \\
Liste ordonnée & O(n) & \textcolor{defvert}{O(1)} & \textcolor{defvert}{O(1)} \\
\rowcolor{defvert!20}
\textbf{Tas} & \textcolor{algoviolet}{O($\log n$)} & \textcolor{defvert}{O(1)} & \textcolor{algoviolet}{O($\log n$)} \\
\hline
\end{tabular}}
\end{center}

\begin{important}
\textbf{Conclusion :} Le tas est la structure optimale pour les files de priorité!
\end{important}

\subsection{TAD Famille d'Ensembles}

\begin{definition}
Soit X un ensemble muni d'une relation d'ordre. On appelle famille les sous-ensembles de X.

\textbf{Interface :} l'insertion d'un ensemble, la suppression d'un ensemble, la vérification d'appartenance d'un ensemble à la famille.
\end{definition}

\begin{notecours}
\textbf{Points importants (notes manuscrites) :}
\begin{itemize}
    \item English "round set", notion d'eun alphabet un ensemble de base = famille ?
    \item Manipuier des sous ensembles
    \item ex {2,4,8} = {4,2,8} (mais interepreter {2,4,8})
    \item Pas d'ordre mais l'ensemble de base a une relation d'ordre
    \item Ex les mots de taille 10 on a 10$^{\wedge}$26 mots possible
    \item La fonction appartenance fait un scan et verifie si= ou pas
    \item C de l ordre de O(taille de famille) on veut ameliorer ca
    \item Et la solution d amelioration c 'est mettre un ordre les sous ensembles sont ordonnes (a,b) n est pas dans la famille
    \item Car non entourer on peut avir des feilles comme des noeud internes
    \item La complexites est devenu de l ordre du taille du l'ensembe debase
    \item Au fur et amesoit on avance et on laisse tombe ce qui est avant donc au pire des cas pour l appartenance il faut tout scaner
    \item Ex :pour ajouter (bde) on fait que avancer pour suppirmer il faut descendre trouver le noeud si cette noeud est une feilles supp jusque j ai plus de feille?
    \item Si c un noued interne c de l ordre de O(taille de l 'alphabet)
\end{itemize}
\end{notecours}

\begin{notecours}
\textbf{Notes sur l'arbre lexicographique (notes manuscrites) :}
\begin{itemize}
    \item \textbf{Pas que un mot plutot un ensemble c plus general}
    \item Table hash --> tableau(calcul) --> tree map --> se balader ds l'arbre c'estu un fonction de hashage aussi a la place des calcules
    \item Un ensemble --> fonction de hashage (calcul) --> @adress (pointeur ds a memoire general) ou est la case
    \item Question a resoudre, un objet e si il appartient a une famille en memoire on utilse pas un tableau ni une liste mais un treemap mapping --> table de hashage
    \item Element <->code<->adress memoire pour 2 ensemble 2 adress different normalment sinon ca s'appele hashage par shainge (avec collision, pas trop abordante, une liste est utiliser)
    \item La complexite de parcourrir l alphabet dans l'arbre
    \item La complexite den map java va dependre de plein d autre truc pour gerer les colision, je chaine ou appeler une autre fonction de hashage
    \item On voit que quand une famille est dens les methodes sont relativement proches
    \item Le mappin grossi norment quand la famille grossi et l'arbre grossi mais de facon maitriser l arbre est plus efficace pour les dense
\end{itemize}
\end{notecours}

\subsection{TAD Ensembles Disjoints (Union-Find)}

\begin{definition}
Un TAD \textbf{ensembles disjoints} gère une famille de sous-ensembles deux à deux disjoints d'un ensemble X, formant une partition.

\textbf{Rappel :} Une partition de E est un ensemble de parties non vides, deux à deux disjointes, dont la réunion est E.
\end{definition}

\begin{notecours}
\textbf{Notes manuscrites sur Union-Find :}
\begin{itemize}
    \item \textbf{E.union(c1,c2) --> O(n)}
    \item \textbf{E.trouverClass(e) --> o(1)}
    \item Si je veux fusionner C2 et C4 je parcoure la table et quand je vois un 4 je le remplace par un 2
    \item Un tas implementer par un table mais juste pour relier le pere d une class par le representant d une autre class, ex pour fusioner (2,4) pere de 4 = 4 devient 2
    \item Au pire des cas (max (hauteur de l arbre))
    \item La complexite depend de quelle class j apporte en fils de l autre, il faut fusioner la plus petit en hauteur je prend ca facien et je le met en fils du plus grande
    \item Je prend la class avec hauteur min et je la branche a l arbre de hauteur max donc on doit savoir le hauteur de l arbre
    \item La hauteur est borne par log(nbr de class)
\end{itemize}
\end{notecours}

\subsubsection{Optimisations}

\begin{important}
\textbf{1. Union par rang :} Attacher l'arbre de plus petite hauteur sous l'arbre de plus grande hauteur.

\textbf{2. Compression de chemin :} Faire pointer tous les nœuds traversés directement vers la racine.
\end{important}

\begin{exemple}
\textbf{Complexité avec optimisations :}

Avec union par rang ET compression de chemin :
\begin{itemize}
    \item m opérations sur n éléments : \textcolor{defvert}{O(m $\cdot$ $\alpha$(n))}
    \item $\alpha$(n) : fonction inverse d'Ackermann (quasi-constante)
    \item En pratique : \textcolor{defvert}{O(1)} amorti par opération
\end{itemize}
\end{exemple}

\newpage
\section{Résumé et Comparaisons}

\subsection{Structures de données}

\begin{center}
\colorbox{fondvert!20}{
\begin{tabular}{|l|p{8cm}|}
\hline
\rowcolor{defvert!30}
\textbf{Structure} & \textbf{Caractéristiques} \\
\hline
Pile & \textcolor{importantrouge}{LIFO}, toutes opérations en \textcolor{defvert}{O(1)} \\
\hline
File & \textcolor{importantrouge}{FIFO}, toutes opérations en \textcolor{defvert}{O(1)} \\
\hline
ABR & Recherche/Insertion/Suppression en \textcolor{algoviolet}{O(h)}. Parcours ordonné efficace \\
\hline
Tas & Accès au max en \textcolor{defvert}{O(1)}, insertion/extraction en \textcolor{algoviolet}{O($\log n$)}. Idéal pour files de priorité \\
\hline
\end{tabular}}
\end{center}

\subsection{Types abstraits de données}

\begin{center}
\colorbox{fondbleu}{
\begin{tabular}{|l|l|p{4.5cm}|}
\hline
\rowcolor{titrebleu!30}
\textbf{TAD} & \textbf{Opérations} & \textbf{Implémentation} \\
\hline
Dictionnaire & insert, search, delete & ABR, table de hachage \\
\hline
Ensemble dynamique & + min, max, pred, succ & Liste chaînée ordonnée \\
\hline
File de priorité & insert, max, extractMax & \textcolor{defvert}{Tas (optimal)} \\
\hline
Famille d'ensembles & insert, delete, member & Arbre lexicographique \\
\hline
Ensembles disjoints & find, union & Forêt d'arbres + optim \\
\hline
\end{tabular}}
\end{center}

\begin{notecours}
\textbf{Question type importante (notes manuscrites) :}

Rappeler que ce que un dictionnaire, implémenter cette structure avec un tas.
\end{notecours}

\newpage
\section{Conclusion}

\begin{important}
Ce cours a présenté les structures de données fondamentales et les types abstraits de données essentiels en informatique. La distinction entre :

\begin{itemize}
    \item \textcolor{algoviolet}{Spécification (TAD)} : définit \textit{quoi} faire
    \item \textcolor{defvert}{Implémentation (structure)} : définit \textit{comment} le faire
\end{itemize}

est cruciale pour concevoir des programmes efficaces et maintenables.
\end{important}

\subsection{Points clés à retenir}

\begin{enumerate}
    \item \textcolor{titrebleu}{Choisir} la structure adaptée selon les opérations nécessaires
    \item \textcolor{titrebleu}{Analyser} les complexités pour optimiser les performances
    \item \textcolor{titrebleu}{Utiliser} les implémentations standards (bibliothèques) quand possible
    \item \textcolor{titrebleu}{Comprendre} les compromis temps/espace
    \item \textcolor{titrebleu}{Maîtriser} la distinction entre interface et implémentation
\end{enumerate}

\subsection{Pour aller plus loin}

\begin{exemple}
\begin{itemize}
    \item Arbres équilibrés (AVL, Rouge-Noir)
    \item Tables de hachage et fonctions de hachage
    \item Graphes et algorithmes associés
    \item Structures de données persistantes
    \item Structures probabilistes (Bloom filter, Skip list)
\end{itemize}
\end{exemple}

\end{document}
